%============================================================ 6. 결과 (압축판)
%  지휘관(LLM) 계층은 본 연구의 기여가 아니라 사용한 구성요소이므로, §6.2 는 표 두 개로
%  줄이고 도판을 두지 않는다. 도판은 기동 계층·체계 전체의 주장에만 쓴다.
\section{결과}
\label{sec:results}

\subsection{기동 계층의 기여}
\label{sec:res_g1}

\subsubsection{휴리스틱 대 점수맵 정책}
\label{sec:res_maneuver}

언어모델 없이 \ref{sec:heuristic}절의 기하 휴리스틱과 \ref{sec:policy}절의 점수맵 정책만
비교한다. 배정은 둘 다 휴리스틱이므로 차이는 기동 방식뿐이다(Table~\ref{tab:maneuver}). 3개
포메이션 통합 기준으로 포획률은 $0.844 \to 0.886$($\dz = 0.37$)으로 $95\,\%$ 수준에서
유의하나, 다중비교 보정을 걸면 구간이 0을 포함한다(\ref{sec:stats}절). 기동 계층이 바꾸는 것은 포획률이 아니라 자원과
기동이다.

\inputshortcap{기동 계층(U-Net 점수맵 + GRPO)의 단독 기여}{tables/tab_maneuver}

\begin{figure}[!tbp]
\centering
\includegraphics[width=\columnwidth]{fig15.png}
\caption{지표 프로파일: a 전체, b 집중, c 양동, d 파상}
\label{fig:radar_maneuver}
\end{figure}

같은 비교를 10개 지표 전부로 넓히면(Fig.~\ref{fig:radar_maneuver}, 각 축은 8개 조건의
범위로 0--1 정규화) U-Net 은 포획당 이동거리 축에서만 안쪽이다. 자리를 골라
그물벽을 놓는 대신 더 멀리 움직이기 때문이다. 충돌률·전개 시 적 거리 축은 겹치고, 나머지 일곱 축에서
바깥에 있다.

포메이션별로 보면 기여가 고르지 않다. 파상에서 $+0.157$($\dz = 0.48$)로 가장 크고, 집중에서는
$0.000$, 양동에서는 $-0.033$ 으로 오히려 소폭 낮다(유의하지 않아 악화로 읽지는 않는다).
\ref{sec:scenario}절의 포메이션 구조가 이 대비를 설명한다. 파상에서는 앞 단을 막은 그물벽의
위치가 뒤 단의 차단을 좌우하므로 해면 전체를 평가하는 점수맵이 직접 효과를 내고, 양동에서는 무리 수와 방어정 수가 같아 결정적인 것이 배정이라 배정이
어긋나면 기동이 보정하지 못한다. 집중은 휴리스틱만으로 포획률 0.997에 도달해 개선
여지가 남지 않는 구간이다(\ref{sec:res_ceiling}절).

\subsubsection{갱신 규칙: 그룹 상대 최적화 대 MAPPO}
\label{sec:res_algo}

점수맵 정책의 이득이 픽셀 지목이라는 표현에서 오는지 가치망 없는 갱신 규칙에서 오는지를
가르기 위해, 같은 액터·행동공간·보상·예산 위에서 갱신 규칙만 MAPPO 로 바꾼 정책 세 개를 같은
프로토콜로 평가하였다(Table~\ref{tab:algo}).

\inputshortcap{같은 픽셀 행동공간 위의 학습 알고리즘 비교}{tables/tab_algo}

포획률은 갈리지 않는다. MAPPO $0.859$ 와 GRPO $0.886$ 의 차이는 유의하지 않고, 포메이션별
방향도 파상 개선·양동 저하로 같다. 갈리는 것은 자원과 조기 제압이다. MAPPO 의 프로파일은
포획률·돌파 축을 빼면 휴리스틱 위에 그대로 겹치는 반면, GRPO 는 포획당 그물을
$0.737 \to 0.568$ 로 줄이고($\dz = -1.19$), 포획 이격거리를 $154\unit{m}$ 늘리며($\dz = +1.21$),
평균 포획 시각과 전개 시 적 거리도 앞당긴다. 네 지표 모두 다중비교 보정에서도 유의하다.
MAPPO 가 이기는 축은 포획당 이동거리 하나다.

같은 표현을 공유하는 한 갱신 규칙은 포획률을 가르지 못하고 그물 절약과 조기 제압만
좌우한다. 휴리스틱 후보를 기준으로 삼는 그룹 상대 이득(\ref{sec:training}절)은 휴리스틱을
넘어서는 선택에만 양의 이득을 주어 그 자원 사용 패턴에서 벗어나게 하는 반면, 팀 가치로
학습하는 MAPPO 는 어느 배가 그물을 아꼈는지를 개별 배에 귀속시킬 신호가 약하다.

\subsection{지휘관 계층의 기여}
\label{sec:res_g2}
\label{sec:res_main}
\label{sec:res_threshold}
\label{sec:res_ceiling}
\label{sec:res_mechanism}
\label{sec:res_latency}

지휘관 계층은 가져다 쓴 구성요소이므로 모델 사이의 순위가 아니라 배치에 필요한 것만 본다.
기동 계층을 휴리스틱으로 고정해 지휘관 축만 분리하면 포획률은 모델 능력을 따라 단조로 오르되,
능력이 낮은 쪽은 휴리스틱 배정 기준선에 미치지 못한다(Table~\ref{tab:main},
Fig.~\ref{fig:condbars}). 개선 폭은 포메이션에 따라 갈려, 휴리스틱만으로 이미 포화한 집중에는
남는 여지가 없고 여유가 없는 양동과 파상에서 크다. 성능을 가르는 것은 배정
커버리지가 아니다(Table~\ref{tab:planquality}). 가장 낮은 성능의 모델이 커버리지는 가장
높았으므로 막아야 할 무리를 빠뜨려서 지는 것이 아니며, 성능과 맞물린 양은 재계획 간 담당
변경(churn)과 경로 교차였다. 재계획마다 방어정의 약 3분의 1이 담당을 바꾸면 그때까지 전개하던
그물벽과 이동이 회수되지 않고, 두 방어정의 진로가 서로를 가로지르는 배정은 충돌 위험을 높인다.
지연도 배치 조건을 가른다. 중형 모델은 응답이 간헐적으로 결정 주기를 넘겨 동기 구조라면 그동안
제어 루프가 멈추므로, \ref{sec:arch}절은 전략 계층만 비동기로 떼어 둔다.

\inputshortcap{조건별 포획률(평균, 95\,\% CI, baseline 대비 차이)}{tables/tab_main}

\begin{figure}[!tbp]
\centering
\includegraphics[width=\columnwidth]{fig18.png}
\caption{포메이션별·조건별 포획률(오차막대 95\,\% CI)}
\label{fig:condbars}
\end{figure}

\inputshortcap{지휘관 계획 품질과 성능}{tables/tab_planquality}

\subsection{제안 체계 전체}
\label{sec:res_g3}

\subsubsection{두 계층의 상호작용}
\label{sec:res_interaction}

식~\eqref{eq:interaction}의 상호작용을 10개 지표 전부에 대해 계산하였다(Table~\ref{tab:interaction}).

\inputshortcap{$2\times2$ 상호작용 분해(Gemini 지휘관)}{tables/tab_interaction}

10개 지표 중 9개에서 상호작용이 유의하지 않고, 다중비교 보정 후에도 그대로 9개다. 두 계층의 기여는 상승이 아니라 가산적이다. 예외인 평균 포획 시각은
유의한 상쇄인데, 두 계층 모두 제압을 앞당기는 쪽으로 작용하므로 같은 포획을 두 번 앞당길
수는 없는 수확체감이다. 가산성이 두 계층을 따로 써도 된다는 뜻은 아니다. 양동에서는 기동
계층 단독($-0.033$)도 전략 계층 단독($+0.070$)도 유의하지 않은데 결합은 $+0.090$ 으로
유의해진다. 무리 수와 방어정 수가 같아 배정과 기동이 둘 다 맞아야 막히기 때문이다.

\subsubsection{전 지표 효과크기}
\label{sec:res_effectsize}

제안 체계와 기준선의 차이를 10개 지표에 걸쳐 효과크기로 잰다(Table~\ref{tab:effectsize},
Fig.~\ref{fig:forest}).

\inputshortcap{전 지표 효과크기($|\dz|$ 내림차순)}{tables/tab_effectsize}

\begin{figure}[!tbp]
\centering
\includegraphics[width=\columnwidth]{fig23.png}
\caption{전 지표 효과크기 $\dz$(옅은 선 보정 구간)}
\label{fig:forest}
\end{figure}

개선은 포획률이 아니라 기동 비용과 자원 효율에 몰려 있다. 총 선회량은 $1625$ 에서 $868$ 로
47\,\% 줄어 $\dz = -3.23$, 포획당 그물은 45\,\% 줄어 $-2.09$, 그물 걸림은 73\,\% 줄어
$-1.19$, 포획 이격거리는 $223\unit{m}$ 늘어 $+1.81$ 이다. 반면 포획률은 $\dz = 0.44$ 로
10개 지표 중 최하위권이며 총 선회량의 7분의 1 수준이다.

다중비교 보정에서는 10개 중 8개가 그대로 유의하고, 포획률과 돌파 수만
유의성을 잃는다.
